\lecture{10}{24. September 2025}{Internal Flows: Fully Developed Laminar Flows}

\begin{problem}{8.42}
  In the laminar flow of an oil of viscosity \qty{1}{\Pa.\s}, the velocity at the center of a \qty{0,3}{\m} pipe is \qty{4,5}{\m\per\s} and the velocity distribution is parabolic.
  Calculate the shear stress at the pipe wall and within the fluid \qty{75}{\mm} from the pipe wall.
\end{problem}

\begin{assumptions}
  \item Steady and incompressible flow
\end{assumptions}

\begin{derivation}
  As the velocity profile is parabolic the flow is fully developed.
  The general formula for the velocity distribution of a fully developed flow in a pipe is
  \[ 
  u = - \frac{R^2}{4 \mu} \left( \frac{\mathrm{d}p}{\mathrm{d}x}  \right) \left( 1 - \left( \frac{r}{R} \right)^2 \right)
  .\]
  In this case the diameter is $D = \qty{0,3}{\m}$ meaning the radius is $R = \frac{D}{2} = \qty{0,15}{\m}$, the flow speed us $u = \qty{4,5}{\m\per\s} $ and the viscosity is $\qty{1}{\Pa.\s}$.
  At the center of the pipe $r = 0$ we have
  \[ 
  u = - \frac{\left( \qty{0,15}{\m}  \right)^2}{4 \cdot \qty{1}{\Pa.\s}} \left( \frac{\mathrm{d}p}{\mathrm{d}x}  \right)
  .\]
  Here we can solve for $\frac{\mathrm{d}p}{\mathrm{d}x} $ to get
  \begin{align*}
  \frac{\mathrm{d}p}{\mathrm{d}x} &= - -\frac{\qty{4,5}{\m\per\s} \cdot 4 \cdot \qty{1}{\Pa.\s}}{\left( \qty{0,15}{\m}  \right)^2} \\
      &= \qty{-800}{\Pa\per\m}
  .\end{align*}
  The shear stress distribution in a pipe is
  \[ 
    \tau_{rx} = \frac{r}{2} \left( \frac{\mathrm{d}p}{\mathrm{d}x}  \right)
  .\]
  At the pipe wall this is
  \[ 
  \tau_{\mathrm{wall}} = \frac{R}{2} \cdot \frac{\mathrm{d}p}{\mathrm{d}x} = \frac{\qty{0,15}{\m} }{2} \cdot \qty{800}{\Pa\per\m} = \qty{60}{\Pa} 
  .\]
  \qty{75}{\mm} from the pipe wall this is
  \[ 
  \tau_{\qty{75}{\mm}} = \frac{\left( \qty{0,15}{\m} - \qty{75}{\mm} \right)}{2} \cdot \qty{800}{\Pa\per\m} = \qty{30}{\Pa} 
  .\]
\end{derivation}

\finalresult{
\begin{aligned}
  \tau_{\mathrm{wall}} &= \qty{60}{\Pa} \\
  \tau_{\qty{75}{\mm} } &= \qty{30}{\Pa}
.\end{aligned}}


\begin{problem}{8.13}
  When a horizontal laminar flow occurs between two parallel plates of infinite extent \qty{0,3}{\m} apart, the velocity at the midpoint between the plates is \qty{2,7}{\m\per\s}. Calculate (a) the flow rate through a section \qty{0,9}{\m} wide, (b) the velocity gradient at the surface of the plate, (c) the wall shearing stress if the fluid has viscosity \qty{1,44}{\Pa.\s}, (d) the pressure drop in each \qty{30}{\m} along the flow. 
\end{problem}

\begin{assumptions}
  \item Steady, incompressible, and fully-established flow
\end{assumptions}

\begin{derivation}
  In this case we must use the velocity distribution for flow between infinite plates,
  \[ 
  u = \frac{a^2}{2\mu} \frac{\mathrm{d}p}{\mathrm{d}x} \left( \left( \frac{y}{a} \right)^2 - \frac{y}{a} \right)
  .\]
  In this particular case $a = \qty{0,3}{\m}$.
  For the velocity at the midpoint ($y = \qty{0,15}{\m}$) we get
  \[ 
  V_c = \frac{\left( \qty{0,3}{\m}  \right)^2}{2 \cdot \mu} \frac{\mathrm{d}p}{\mathrm{d}x}  \left( \left( \frac{\qty{0,15}{\m} }{\qty{0,3}{\m} } \right)^2 - \left( \frac{\qty{0,15}{\m} }{\qty{0,30}{\m} } \right) \right) = \qty{2,7}{\m\per\s} 
  .\]
  And therefore the pressure gradient can be found as
  \[ 
    \frac{1}{2\mu}\frac{\mathrm{d}p}{\mathrm{d}x} = \frac{\qty{2,7}{\m\per\s}}{\left( \qty{0,3}{\m}  \right)^2 \cdot \left( \left( \frac{\qty{0,15}{\m} }{\qty{0,3}{\m} } \right)^2 - \frac{\qty{0,15}{\m} }{\qty{0,30}{\m} } \right)} = \qty{-120}{\per\m\per\s} 
  .\]
  The average velocity is given by
  \[ 
  \overline{V} = - \frac{1}{12\mu} \frac{\mathrm{d}p}{\mathrm{d}x} a^2
  .\]
  Inputting known values yields
  \[ 
  \overline{V} = \qty{120}{\per\m\per\s} \cdot \frac{1}{6} \cdot \left( \qty{0,3}{\m}  \right)^2 = \qty{1,8}{\m\per\s} 
  .\]
  Which gives a volumetric flow rate of
  \[ 
  Q = A \cdot \overline{V} = \qty{0,3}{\m} \cdot \qty{0,9}{\m} \cdot \qty{1,8}{\m\per\s}  = \qty{0,486}{\cubic\m\per\s} 
  .\]
  
  The velocity gradient is
  \[ 
  \frac{\mathrm{d}u}{\mathrm{d}y} = \frac{a^2}{2\mu} \frac{\mathrm{d}p}{\mathrm{d}x} \left( \frac{2y}{a^2} - \frac{1}{a} \right) = \left( \qty{0,3}{\m}  \right)^2 \cdot \left( \qty{-120}{\per\m\per\s} \right) \cdot \left( 0 - \frac{1}{\qty{0,30}{\m} } \right) = \qty{36}{\per\s} 
  .\]

  The shear stress at the wall is
  \[ 
  \tau_{yx} = \mu \frac{\mathrm{d}u}{\mathrm{d}y} = \qty{1,44}{\Pa.\s} \cdot \qty{36}{\per\s} = \qty{51,84}{\Pa} 
  .\]

  For the pressure drop per \qty{30}{\m} we know that
  \begin{align*}
    \frac{1}{2\mu} \frac{\mathrm{d}p}{\mathrm{d}x} &= \qty{-120}{\per\m\per\s}  \\
    \frac{\mathrm{d}p}{\mathrm{d}x} &= \qty{-120}{\per\m\per\s} \cdot 2 \cdot \qty{1,44}{\Pa.\s}  \\
    &= \qty{-345,6}{\Pa\per\m}  \\
  .\end{align*}
  Now entering the length $\Delta x = \qty{30}{\m} $ we get
  \[ 
  \Delta p_{\qty{30}{\m} } = \frac{\mathrm{d}p}{\mathrm{d}x} \cdot (-\Delta x) = \qty{-345,6}{\Pa\per\m} \cdot \qty{-30}{\m} = \qty{10,368}{\kPa} 
  .\]
\end{derivation}

\finalresult{
\begin{aligned}
  Q_{\qty{0,9}{\m} } &= \qty{0,486}{\cubic\m\per\s}  \\
  \left( \frac{\mathrm{d}u}{\mathrm{d}y}  \right)_{\mathrm{wall}} &= \qty{36}{\per\s} \\
  \tau_{\text{wall}} &= \qty{51,8}{\Pa}  \\
  \Delta p_{\qty{30}{\m} } &= \qty{10,4}{\kPa} 
.\end{aligned}
}

\begin{problemwithimage}{p10_1}{8.6}
  A fully developed and laminar flow of oil occurs between parallel plates.
  The pressure gradient creating the flow is \qty{1,25}{\kPa\per\m} of length and the channel half-width is \qty{1,5}{\mm}.
  Calculate the magnitude and length direction of the wall shear stress at the upper plate surface.
  Find the volume flow rate through the channel.
  The viscosity is \qty{0,50}{\N.\s\per\square\m}. 
\end{problemwithimage}

\begin{derivation}
  The basic velocity distribution is
  \[ 
  u = - \frac{a^2}{2\mu} \frac{\mathrm{d}p}{\mathrm{d}x} \left( 1 - \left( \frac{y}{a} \right)^2 \right)
  .\]
  This corresponds to a velocity gradient of
  \[ 
  \frac{\mathrm{d}u}{\mathrm{d}y} = - \frac{a^2}{2\mu} \frac{\mathrm{d}p}{\mathrm{d}x} \left( \frac{2y}{a^2} \right)
  .\]
  Now the shear stress distribution can be found as
  \[ 
  \tau_{yx} = \mu \frac{\mathrm{d}u}{\mathrm{d}y} = - \frac{a^2}{2} \frac{\mathrm{d}p}{\mathrm{d}x} \left( - \frac{2y}{a^2} \right) = -y \frac{\mathrm{d}p}{\mathrm{d}x} 
  .\]
  At the outer edge where $y = h$ this becomes
  \[ 
  \tau_{yx} = -h \frac{\mathrm{d}p}{\mathrm{d}x} = -\qty{1,5}{\mm} \cdot \qty{1,25}{\kPa\per\m} = \qty{-1,88}{\Pa}
  .\]
  The volume flow rate is
  \[ 
  Q = \int u \, \mathrm{d}A = \int_{-h}^{h} u \cdot b \, \mathrm{d}y = -\frac{b h^2}{2\mu} \frac{\mathrm{d}p}{\mathrm{d}x} \cdot \int_{-h}^{h} 1 - \left( \frac{y}{h} \right)^2 \, \mathrm{d}y  = - \frac{2 h^3 \cdot b}{3\mu} \frac{\mathrm{d}p}{\mathrm{d}x} 
  .\] 
  Now we can solve for the flow rate per unit depth $\frac{Q}{b}$ to get
  \[ 
  \frac{Q}{b} = - \frac{2}{3} \frac{h^3}{\mu} \frac{\mathrm{d}p}{\mathrm{d}x} = - \frac{2}{3} \cdot \left( \qty{1,5}{\mm}  \right)^3 \cdot \frac{1}{\qty{0,50}{\N\s\per\square\m}} \cdot \qty{1,25}{\kPa\per\m} = \qty{-5,6}{\square\mm\per\s} 
  .\]
\end{derivation}

\finalresult{\frac{Q}{b} = \qty{-5,6e-6}{\square\m\per\s}}
